\documentclass[12pt, a4paper]{article}

% ── PACKAGES ──────────────────────────────────────────────────
\usepackage[utf8]{inputenc}
\usepackage[T1]{fontenc}
\usepackage{geometry}
\geometry{margin=2.5cm}
\usepackage{hyperref}
\hypersetup{
    colorlinks=true,
    linkcolor=blue,
    urlcolor=blue,
    citecolor=blue,
    pdftitle={EZH2 Inhibition Combined with Anti-HER2 Therapy
              for the HER2-Deep Fraction: CDH3-High, AR-Low,
              EZH2 +118 Percent},
    pdfauthor={Eric Robert Lawson},
    pdfkeywords={EZH2, EZH2 inhibitor, tazemetostat,
                 HER2, HER2-enriched, anti-HER2,
                 trastuzumab, pertuzumab, CDH3, P-cadherin,
                 androgen receptor, AR, de-differentiation,
                 epigenetic lock, Waddington landscape,
                 attractor geometry, breast cancer,
                 subpopulation, resistance, GSE176078,
                 OrganismCore, HER2 deep fraction,
                 undertreated, cell fate plasticity}
}
\usepackage{amsmath}
\usepackage{booktabs}
\usepackage{longtable}
\usepackage{array}
\usepackage{parskip}
\usepackage{microtype}
\usepackage{authblk}
\usepackage{titlesec}
\usepackage{fancyhdr}
\usepackage{xcolor}
\usepackage{float}
\usepackage{enumitem}
\usepackage{setspace}

% ── COLOURS ───────────────────────────────────────────────────
\definecolor{repobox}{RGB}{240,247,255}
\definecolor{findingbox}{RGB}{245,255,245}
\definecolor{novelbox}{RGB}{255,248,220}
\definecolor{mechanismbox}{RGB}{250,245,255}
\definecolor{actionbox}{RGB}{230,255,230}
\definecolor{urgentbox}{RGB}{255,235,235}
\definecolor{convergentbox}{RGB}{255,250,235}
\definecolor{contrastbox}{RGB}{248,240,255}

% ── HEADER / FOOTER ───────────────────────────────────────────
\pagestyle{fancy}
\fancyhf{}
\lhead{\footnotesize CS-LIT-21: EZH2i + Anti-HER2\\
       \footnotesize for the HER2-Deep Fraction}
\rhead{\footnotesize OrganismCore \textbar{} 2026-03-06}
\rfoot{\small Page \thepage}
\renewcommand{\headrulewidth}{0.4pt}
\setlength{\headheight}{24pt}
\addtolength{\topmargin}{-10pt}

% ── SECTION FORMATTING ────────────────────────────────────────
\titleformat{\section}
  {\large\bfseries}{\thesection}{1em}{}[\titlerule]
\titleformat{\subsection}
  {\normalsize\bfseries}{\thesubsection}{1em}{}
\titleformat{\subsubsection}
  {\small\bfseries}{\thesubsubsection}{1em}{}

% ══════════════════════════════════════════════════════════════
% TITLE
% ══════════════════════════════════════════════════════════════
\title{
    \vspace{-1em}
    \textbf{EZH2 Inhibition Combined with Anti-HER2
    Therapy for the HER2-Deep Fraction:}\\[0.5em]
    \large CDH3-High, AR-Low, EZH2 $+118\%$ ---
    A De-Differentiated Subpopulation Within
    HER2-Enriched Breast Cancer That Is Currently
    Undertreated by Standard Anti-HER2 Regimens\\[0.4em]
    \normalsize \textit{OrganismCore Framework ---
    Document CS-LIT-21}\\[0.3em]
    \small Derived from Waddington Attractor Geometry
    of 19,542 Single Cancer Cells (GSE176078)\\[0.3em]
    \small Part of the BRCA Cross-Subtype Analysis
    (BRCA-S8h, 2026-03-05)
}

\author[1]{Eric Robert Lawson}
\affil[1]{
    OrganismCore\\
    \href{mailto:OrganismCore@proton.me}
         {OrganismCore@proton.me}\\
    ORCID:
    \href{https://orcid.org/0009-0002-0414-6544}
         {0009-0002-0414-6544}
}

\date{March 6, 2026}

% ══════════════════════════════════════════════════════════════
\begin{document}
% ══════════════════════════════════════════════════════════════

\maketitle
\thispagestyle{fancy}

% ── REPOSITORY POINTER ────────────────────────────────────────
\begin{center}
\colorbox{repobox}{
\begin{minipage}{0.88\textwidth}
\vspace{0.5em}
\centering
\textbf{Full computational record --- scripts, data
caches, reasoning artifacts, and raw logs:}\\[0.4em]
\href{https://github.com/Eric-Robert-Lawson/attractor-oncology}
{\texttt{https://github.com/Eric-Robert-Lawson/attractor-oncology}}
\\[0.3em]
\small Part of the OrganismCore BRCA Cross-Subtype
Analysis (BRCA-S8h, 2026-03-05).\\
All predictions were locked from attractor geometry
\textbf{before} any literature review was performed.
\vspace{0.5em}
\end{minipage}}
\end{center}

\vspace{1em}

% ── FINDING BOX ───────────────────────────────────────────────
\begin{center}
\colorbox{findingbox}{
\begin{minipage}{0.88\textwidth}
\vspace{0.5em}
\textbf{THE FINDING, STATED PRECISELY:}\\[0.4em]
Within the HER2-enriched subtype in GSE176078,
Waddington attractor geometry identifies a
distinct de-differentiated subpopulation
characterised by three co-occurring features:

\begin{itemize}[noitemsep]
    \item \textbf{CDH3-high} (P-cadherin
    overexpression: marker of de-differentiation
    and mesenchymal/stem-like identity)
    \item \textbf{AR-low} (androgen receptor
    suppressed: loss of luminal identity program)
    \item \textbf{EZH2 $+118\%$} (EZH2 elevated
    118\% above the HER2-enriched subtype mean)
\end{itemize}

This population exists \textbf{within} the
HER2-enriched subtype. It is not a separate
subtype. It is the deeply locked fraction of
HER2+ disease --- the cells that have partially
escaped HER2-dependent identity and descended
into an EZH2-driven epigenetic lock. Standard
anti-HER2 therapy (trastuzumab/pertuzumab)
targets the HER2 receptor. It does not address
the EZH2-driven de-differentiation that
characterises this fraction.
\vspace{0.5em}
\end{minipage}}
\end{center}

\vspace{1em}

% ── CONVERGENCE BOX ───────────────────────────────────────────
\begin{center}
\colorbox{convergentbox}{
\begin{minipage}{0.88\textwidth}
\vspace{0.5em}
\textbf{INDEPENDENT CONVERGENCE:}\\[0.3em]
Muller et al.\ (EMBO Reports, 2026) independently
demonstrated that EZH2 drives HER2+ breast cancer
progression through cell-fate plasticity and
de-differentiation. Their finding that EZH2
inhibition shifts HER2+ tumor cells toward a
luminal phenotype is the mechanistic confirmation
of the attractor geometry prediction made here.
The two findings are orthogonal in method and
convergent in conclusion.
\vspace{0.5em}
\end{minipage}}
\end{center}

\vspace{1em}

% ══════════════════════════════════════════════════════════════
\begin{abstract}
% ══════════════════════════════════════════════════════════════
Standard HER2-enriched breast cancer is treated
with HER2-directed monoclonal antibodies
(trastuzumab, pertuzumab) combined with
chemotherapy. This approach achieves high
pathological complete response rates in many
patients. However, a subpopulation within
HER2-enriched disease shows primary or acquired
resistance to anti-HER2 therapy that is not
explained by HER2 pathway genetics alone.
We report the identification by Waddington
attractor geometry, applied to 19,542 single
cancer cells (GSE176078), of a de-differentiated
subpopulation within HER2-enriched breast cancer
defined by CDH3 overexpression, AR suppression,
and EZH2 elevation ($+118\%$ above the
HER2-enriched mean). This population, termed the
\textbf{HER2-Deep Fraction}, has partially
escaped HER2-dependent luminal identity and
descended into an EZH2-driven epigenetic lock
state. The geometry predicts that anti-HER2
therapy alone is insufficient for this fraction:
the EZH2 lock must be addressed concurrently.
The proposed intervention is EZH2 inhibition
(tazemetostat) combined with standard anti-HER2
therapy, with CDH3 expression as the patient
selection biomarker for the deep fraction.
This finding independently converges with
recent published work demonstrating that EZH2
drives HER2+ breast cancer progression through
cell-fate plasticity (Muller et al., EMBO
Reports, 2026). No published instrument
identifies the HER2-deep fraction by attractor
geometry using CDH3-high, AR-low, and EZH2+118\%
as co-occurring markers. The finding is novel.
\end{abstract}

\tableofcontents
\newpage

% ═══════════════════════════════════════════���══════════════════
\section{Context and Origin}
% ══════════════════════════════════════════════════════════════

This document covers CS-LIT-21 from the
OrganismCore BRCA Cross-Subtype Literature Check
(BRCA-S8h, 2026-03-05):

\begin{quote}
\textbf{CS-LIT-21}: EZH2i + anti-HER2 for the
HER2-deep fraction (CDH3-high, AR-low,
EZH2 $+118\%$)
\end{quote}

The finding was derived from attractor geometry
before any literature review was performed.
The literature check assigned it a
\textbf{NOVEL} verdict, with a note that it is
directly relevant to trastuzumab/pertuzumab
resistance research.

\subsection{Relationship to Other Published
Documents in This Series}

\begin{itemize}
    \item \textbf{CS-LIT-2 (Six Lock Type
    Classification)}\\
    \href{https://doi.org/10.5281/zenodo.18884158}
    {DOI: 10.5281/zenodo.18884158}\\
    HER2-enriched is Lock Type 2 in the
    six-lock taxonomy. The HER2-deep fraction
    represents the sub-population within
    Lock Type 2 that has acquired properties
    of Lock Type 3 (the deep EZH2 lock of TNBC).

    \item \textbf{CS-LIT-1 (FOXA1/EZH2 Ratio)}\\
    \href{https://doi.org/10.5281/zenodo.18883922}
    {DOI: 10.5281/zenodo.18883922}\\
    The FOXA1/EZH2 ratio orders subtypes
    across the full spectrum. The HER2-deep
    fraction is the subpopulation of
    HER2-enriched with the lowest FOXA1/EZH2
    ratio within the subtype --- the cells
    most shifted toward the TNBC end of the
    continuous ordering axis.

    \item \textbf{CS-LIT-6 (AR as Continuous
    Depth Axis)}\\
    \href{https://doi.org/10.5281/zenodo.18891770}
    {DOI: 10.5281/zenodo.18891770}\\
    AR suppression in the HER2-deep fraction
    mirrors the AR-low signature of deeply
    locked TNBC. The same continuous depth
    axis operates across both subtypes.

    \item \textbf{CS-LIT-16 (Tazemetostat
    $\rightarrow$ Fulvestrant in TNBC)}\\
    \href{https://doi.org/10.5281/zenodo.18884003}
    {DOI: 10.5281/zenodo.18884003}\\
    The EZH2 inhibitor logic developed for
    TNBC applies to the HER2-deep fraction
    through the same attractor geometry
    principle. The difference is the
    combination: in HER2-deep, the partner
    is anti-HER2 therapy rather than
    fulvestrant.
\end{itemize}

% ══════════════════════════════════════════════════════════════
\section{The Problem: What Standard Anti-HER2
Therapy Does and Does Not Address}
% ══════════════════════════════════════════════════════════════

\subsection{What Standard Anti-HER2 Therapy
Targets}

Trastuzumab and pertuzumab bind HER2
(ERBB2) extracellularly and inhibit:
\begin{itemize}
    \item HER2 homodimerisation and
    heterodimerisation
    \item Downstream PI3K/AKT and MAPK/ERK
    signalling
    \item Antibody-dependent cellular
    cytotoxicity (ADCC)
\end{itemize}

These mechanisms are effective for tumor
cells whose survival is driven by
HER2 receptor signalling. The assumption
is that HER2-enriched tumors are
predominantly HER2-dependent.

\subsection{What the Attractor Geometry
Reveals}

The Waddington landscape of HER2-enriched
cells in GSE176078 is not homogeneous.
Within the HER2-enriched cell population,
there is a distinct fraction of cells that
have partially disengaged from
HER2-dependent luminal identity and
descended into an EZH2-driven attractor
state. These cells are defined by:

\begin{center}
\begin{tabular}{p{3.5cm}p{5.5cm}p{3.5cm}}
\toprule
\textbf{Marker} &
\textbf{Biological meaning} &
\textbf{Value in deep fraction} \\
\midrule
CDH3 (P-cadherin) &
Marks de-differentiated, mesenchymal-like,
stem-cell-adjacent state. Associated with
loss of epithelial identity and increased
migratory potential. &
High (elevated above HER2 mean) \\[0.5em]
AR &
Continuous readout of luminal identity
program activity (see CS-LIT-6). Low AR
indicates partial loss of luminal identity
and deeper attractor descent. &
Low (suppressed) \\[0.5em]
EZH2 &
Primary driver of the epigenetic lock.
EZH2 $+118\%$ above HER2-enriched mean
indicates active PRC2-mediated epigenetic
silencing in this fraction. &
$+118\%$ above subtype mean \\
\bottomrule
\end{tabular}
\end{center}

\vspace{0.5em}

The co-occurrence of these three features
in the same cells is not coincidental.
It is a geometric signature: these cells
are in a different attractor position
from the majority of HER2-enriched cells.
They have not simply amplified HER2
expression --- they have additionally
activated the EZH2/PRC2 program that
drives epigenetic de-differentiation.

\subsection{Why This Creates a Therapeutic
Gap}

Anti-HER2 therapy targets the HER2 receptor.
It does not:
\begin{itemize}
    \item Inhibit EZH2-driven epigenetic
    silencing
    \item Restore luminal identity programs
    suppressed by PRC2
    \item Address CDH3-mediated mesenchymal
    plasticity
    \item Reduce the EZH2 activity that
    maintains the de-differentiated state
\end{itemize}

The HER2-deep fraction therefore has an
epigenetic survival mechanism that is
orthogonal to HER2 receptor inhibition.
Blocking HER2 alone leaves this mechanism
intact. The prediction is that the HER2-deep
fraction is enriched in anti-HER2
non-responders and in patients who relapse
after initial response.

% ══════════════════════════════════════════════════════════════
\section{The Three-Marker Geometric Signature}
% ══════════════════════════════════════════════════════════════

\subsection{CDH3 as the Primary Identifier
of the Deep Fraction}

CDH3 (P-cadherin) is a cell adhesion molecule
whose overexpression in breast cancer marks:
\begin{itemize}
    \item Loss of E-cadherin-mediated
    epithelial organisation
    \item Acquisition of stem-cell-adjacent
    phenotype
    \item Increased migratory and invasive
    capacity
    \item Resistance to conventional
    cytotoxic agents
\end{itemize}

In the Waddington landscape, CDH3
overexpression within HER2-enriched cells
marks the cells that have descended from the
canonical HER2 luminal attractor toward
a more basal/de-differentiated position.
CDH3 is high specifically in the cells that
have also activated EZH2 --- it is not
elevated uniformly across HER2-enriched
cells.

CDH3-high therefore functions as the primary
clinical identifier of the deep fraction:
it is measurable by IHC and RNA-seq in
standard clinical samples.

\subsection{AR Suppression as Confirmation
of Attractor Depth}

Within HER2-enriched cells, AR suppression
mirrors the same continuous depth axis
established for TNBC in CS-LIT-6
(\href{https://doi.org/10.5281/zenodo.18891770}
{DOI: 10.5281/zenodo.18891770}). The cells
with the lowest AR expression within
HER2-enriched are the same cells with the
highest EZH2 and the highest CDH3. The
three-marker co-occurrence is a single
geometric event, not three independent
associations.

\subsection{\texorpdfstring{EZH2 $+118\%$:
The Magnitude of the Lock Within the
Subtype}{EZH2 +118 Percent: The Magnitude
of the Lock}}

EZH2 is elevated $+118\%$ above the
HER2-enriched subtype mean in the deep
fraction. This is a substantial within-subtype
elevation --- not merely a reflection of
the already-elevated EZH2 characteristic
of HER2-enriched relative to luminal subtypes.
It represents the additional EZH2 activity
that specifically characterises the
de-differentiated subpopulation.

For comparison: EZH2 is elevated $+189\%$
above baseline in TNBC (CS-LIT-4 in BRCA-S8h).
The HER2-deep fraction at $+118\%$ is
intermediate between typical HER2-enriched
and typical TNBC --- consistent with the
geometric interpretation that these cells
are partially descending toward the TNBC
attractor while remaining HER2-amplified.

% ══════════════════════════════════════════════════════════════
\section{Independent Convergence with Published
Literature}
% ══════════════════════════════════════════════════════════════

\begin{center}
\colorbox{convergentbox}{
\begin{minipage}{0.88\textwidth}
\vspace{0.5em}
\textbf{CONVERGENT FINDING:}\\[0.4em]
\textbf{Muller et al., EMBO Reports, 2026:}\\
\textit{EZH2 directs HER2+ breast cancer
progression through cell-fate plasticity.}\\[0.4em]
This study, published independently of the
OrganismCore analysis, demonstrates using
genetically engineered mouse models that
EZH2 is required for HER2+ breast cancer
progression. EZH2 inhibition shifts HER2+
tumor cells toward a luminal, less aggressive
phenotype and sensitises them to endocrine
therapy. The mechanism is cell-fate
plasticity: EZH2 maintains the de-differentiated
state of a fraction of HER2+ tumor cells.\\[0.4em]
\textbf{Why this matters:}\\
The Muller et al. study confirms the
mechanistic basis of the CS-LIT-21 prediction
by an independent method (mouse model
genetics) in an independent laboratory.
The EZH2-driven de-differentiation of
HER2+ cells is now mechanistically established
in the published literature. CS-LIT-21
identifies the human clinical counterpart
of this mechanism using attractor geometry
applied to human single-cell data.
\vspace{0.5em}
\end{minipage}}
\end{center}

\vspace{0.5em}

Additionally, Bao et al.\ (Nature Communications,
2020) demonstrated that EZH2-mediated repression
of PP2A (via PPP2R2B silencing) confers
resistance to HER2-targeted therapy by
sustaining PI3K/AKT/mTOR signalling downstream
of HER2 blockade. This is the molecular
mechanism by which EZH2 activity in the
HER2-deep fraction enables survival despite
anti-HER2 therapy.

A Phase 1b study of the EZH1/2 inhibitor
valemetostat in combination with trastuzumab
deruxtecan (T-DXd) in breast cancer is
ongoing (ASCO 2024 abstract TPS1120),
confirming that the combination strategy
is now entering clinical testing in an
adjacent context.

% ══════════════════════════════════════════════════════��═══════
\section{The Proposed Intervention}
% ══════════════════════════════════════════════════════════════

\subsection{The Combination Logic}

\begin{center}
\colorbox{mechanismbox}{
\begin{minipage}{0.88\textwidth}
\vspace{0.5em}
\textbf{Two orthogonal mechanisms, two
targeted agents:}\\[0.5em]
\textbf{Anti-HER2 therapy} (trastuzumab
and/or pertuzumab):\\
Targets HER2 receptor signalling.
Addresses the HER2-dependent survival
mechanism present in all HER2-enriched
cells, including the deep fraction.\\[0.5em]
\textbf{EZH2 inhibitor} (tazemetostat):\\
Targets the EZH2/PRC2-driven epigenetic
lock. Addresses the de-differentiation
and epigenetic survival mechanism that
is specifically elevated in the deep
fraction. Restores luminal identity
programs suppressed by PRC2 activity.
Reduces CDH3 expression and mesenchymal
plasticity as a downstream consequence
of lock reversal.\\[0.5em]
\textbf{Combined effect}:\\
Anti-HER2 eliminates HER2-dependent
survival; EZH2 inhibition eliminates the
epigenetic survival mechanism that
otherwise allows the deep fraction to
persist and eventually repopulate the
tumor after anti-HER2 therapy.
\vspace{0.5em}
\end{minipage}}
\end{center}

\subsection{Patient Selection Biomarker}

CDH3 expression (by IHC or RNA-seq) is the
proposed patient selection biomarker for
enriching the HER2-deep fraction:

\begin{center}
\begin{tabular}{p{4cm}p{4cm}p{4.5cm}}
\toprule
\textbf{CDH3 Status} &
\textbf{Geometric Position} &
\textbf{Clinical Implication} \\
\midrule
CDH3-high (upper quartile
within HER2-enriched) &
Deep fraction. EZH2 elevated.
AR suppressed. Partially
de-differentiated. &
Primary target for EZH2i +
anti-HER2 combination.
High probability of
insufficient response
to anti-HER2 alone. \\[0.5em]
CDH3-intermediate &
Transitional. Some EZH2
activity but not deep lock. &
May benefit; lower priority
for combination than CDH3-high. \\[0.5em]
CDH3-low (standard
HER2-enriched) &
Shallow. HER2-dependent
survival dominant. &
Standard anti-HER2 therapy
appropriate. EZH2i addition
not mechanistically justified. \\
\bottomrule
\end{tabular}
\end{center}

\subsection{Sequencing Considerations}

Two sequencing strategies are geometrically
motivated:

\begin{enumerate}
    \item \textbf{Concurrent combination:}
    EZH2 inhibitor administered alongside
    anti-HER2 therapy from initiation.
    Rationale: prevents the deep fraction
    from establishing an epigenetically
    locked residual population during
    the period of HER2 blockade.

    \item \textbf{EZH2i priming followed
    by anti-HER2:} EZH2 inhibitor
    administered first to reverse the
    epigenetic lock and restore luminal
    identity, followed by anti-HER2.
    Rationale: analogous to the
    tazemetostat $\rightarrow$ fulvestrant
    sequence in TNBC (CS-LIT-16), where
    lock reversal unmasks a subsequent
    hormonal target. In HER2-deep, lock
    reversal may restore full
    HER2-dependency in previously
    resistant cells, enhancing subsequent
    anti-HER2 efficacy.
\end{enumerate}

% ══════════════════════════════════════════════════════════════
\section{What the Published Literature Does and
Does Not Contain}
% ══════════════════════════════════════════════════════════════

\subsection{What IS in the Published Literature}

\begin{itemize}
    \item EZH2 inhibition sensitises
    HER2+ cells to anti-HER2 therapy
    via PP2A restoration (Bao et al.,
    Nature Communications, 2020).
    \item EZH2 drives HER2+ breast cancer
    progression through de-differentiation
    and cell-fate plasticity (Muller et al.,
    EMBO Reports, 2026).
    \item EZH1/2 inhibitor combined with
    T-DXd is in Phase 1b testing in
    breast cancer (ASCO 2024).
    \item CDH3 (P-cadherin) overexpression
    marks aggressive, de-differentiated
    breast cancer subpopulations.
    \item AR suppression is associated with
    de-differentiation in breast cancer.
\end{itemize}

\subsection{What is NOT in the Published
Literature}

\begin{itemize}
    \item The identification of the
    HER2-deep fraction as a geometric
    entity defined by the co-occurrence
    of CDH3-high, AR-low, and
    EZH2 $+118\%$ within HER2-enriched
    breast cancer.
    \item The use of CDH3 expression as
    a patient selection biomarker for
    EZH2 inhibitor benefit specifically
    within HER2-enriched disease.
    \item The framing of the HER2-deep
    fraction as a partially descended
    attractor state --- cells that are
    HER2-amplified but have additionally
    activated the EZH2-driven lock of
    TNBC.
    \item The specific combination
    proposal (EZH2i + anti-HER2) derived
    from attractor geometry applied to
    human single-cell data (GSE176078).
    \item The EZH2 $+118\%$ quantification
    within HER2-enriched cells as a
    within-subtype elevation distinct
    from the between-subtype elevation.
\end{itemize}

\begin{center}
\colorbox{novelbox}{
\begin{minipage}{0.88\textwidth}
\vspace{0.5em}
\textbf{VERDICT: NOVEL}\\[0.3em]
The identification of the HER2-deep fraction
by attractor geometry (CDH3-high, AR-low,
EZH2 $+118\%$) and the derived combination
proposal (EZH2 inhibitor + anti-HER2) have
no direct published equivalent as of
2026-03-06. The finding independently
converges with Muller et al.\ (2026) on
the mechanism (EZH2-driven
de-differentiation in HER2+ breast cancer)
while providing a distinct patient selection
instrument and combination strategy not
present in the published literature.
\vspace{0.5em}
\end{minipage}}
\end{center}

% ══════════════════════════════════════════════════════════════
\section{Validation Pathway}
% ══════════════════════════════════════════════════════════════

\begin{center}
\colorbox{actionbox}{
\begin{minipage}{0.88\textwidth}
\vspace{0.5em}
\textbf{Proposed analyses:}\\[0.4em]
\textbf{Analysis 1 --- CDH3 subgroup
outcome analysis:}\\
In a HER2-enriched cohort with anti-HER2
treatment and outcome data: stratify by
CDH3 expression (IHC H-score or RNA-seq
quartile). Test whether CDH3-high patients
have worse pCR rates and/or worse PFS
on standard anti-HER2 therapy.
This is the falsifiable prediction from
the geometric model.\\[0.5em]
\textbf{Analysis 2 --- EZH2 / CDH3
co-expression confirmation:}\\
In GSE176078 or an independent single-cell
HER2-enriched dataset: confirm that CDH3-high
and EZH2-high are co-occurring in the same
cells rather than in separate subpopulations.
The geometric claim requires co-occurrence.\\[0.5em]
\textbf{Analysis 3 --- Combination
efficacy in model systems:}\\
In a HER2+ breast cancer cell line with
CDH3 overexpression and EZH2 elevation:
test tazemetostat + trastuzumab combination
vs each agent alone. Measure CDH3 expression,
EZH2 activity, cell viability, and apoptosis.\\[0.5em]
\textbf{Analysis 4 --- IHC cohort
stratification:}\\
Apply FOXA1/EZH2 dual IHC with CDH3 to
an archival HER2-enriched cohort. Quantify
the CDH3-high, EZH2-high subpopulation
fraction and correlate with anti-HER2
response and survival.
\vspace{0.5em}
\end{minipage}}
\end{center}

% ══════════════════════════════════════════════════════════════
\section{What This Document Claims and Does
Not Claim}
% ══════════════════════════════════════════════════════════════

\subsection{What IS Claimed}

\begin{itemize}
    \item Within HER2-enriched cells in
    GSE176078, a de-differentiated
    subpopulation is identifiable by
    the co-occurrence of CDH3-high,
    AR-low, and EZH2 $+118\%$.
    \item This subpopulation occupies a
    distinct attractor position from the
    majority of HER2-enriched cells.
    \item The EZH2 elevation in this
    fraction is within-subtype, not
    merely a reflection of the
    between-subtype EZH2 gradient.
    \item Standard anti-HER2 therapy
    does not address the EZH2-driven
    epigenetic lock mechanism in this
    fraction.
    \item The combination of EZH2i +
    anti-HER2 is mechanistically
    motivated by the geometry.
    \item CDH3 expression is the proposed
    patient selection biomarker for the
    deep fraction.
    \item The finding independently
    converges with Muller et al.\ (2026)
    on EZH2-driven de-differentiation
    in HER2+ breast cancer.
\end{itemize}

\subsection{What is NOT Claimed}

\begin{itemize}
    \item The HER2-deep fraction has not
    been prospectively validated in a
    clinical cohort.
    \item The combination of EZH2i +
    anti-HER2 has not been tested in
    CDH3-selected patients.
    \item CDH3 has not been validated as
    a predictive biomarker for EZH2i
    benefit in HER2-enriched disease.
    \item The finding is from a single
    single-cell dataset and requires
    replication in independent data.
\end{itemize}

% ══════════════════════════════════════════════════════════════
\section{Locked Statement}
% ══════════════════════════════════════════════════════════════

\begin{center}
\colorbox{findingbox}{
\begin{minipage}{0.88\textwidth}
\vspace{0.6em}
\textbf{Stated once, precisely, without
inflation:}\\[0.5em]
Within the HER2-enriched cell population
of GSE176078, Waddington attractor geometry
identifies a de-differentiated subpopulation
defined by CDH3 overexpression, AR suppression,
and EZH2 elevation ($+118\%$ above the
HER2-enriched subtype mean). This population
--- the HER2-Deep Fraction --- has activated
the same EZH2-driven epigenetic lock that
characterises deeply locked TNBC, while
remaining HER2-amplified. Standard anti-HER2
therapy does not address the EZH2 mechanism.
The proposed intervention is EZH2 inhibitor
(tazemetostat) combined with anti-HER2 therapy,
with CDH3 expression as the patient selection
biomarker. This finding independently
converges with Muller et al.\ (EMBO Reports,
2026) on EZH2-driven cell-fate plasticity
in HER2+ breast cancer. The specific
three-marker signature, attractor geometry
derivation, and CDH3-based patient selection
proposal have no direct published equivalent
as of 2026-03-06.
\vspace{0.5em}
\end{minipage}}
\end{center}

% ══════════════════════════════════════════════════════════════
\section{Document Metadata}
% ══════════════════════════════════════════════════════════════

\begin{tabular}{ll}
\toprule
\textbf{Field} & \textbf{Value} \\
\midrule
Document ID & CS-LIT-21 \\
Parent document & BRCA-S8h (2026-03-05) \\
Verdict & NOVEL \\
Date & 2026-03-06 \\
Author & Eric Robert Lawson \\
ORCID & 0009-0002-0414-6544 \\
Contact & OrganismCore@proton.me \\
Repository &
\href{https://github.com/Eric-Robert-Lawson/attractor-oncology}
{attractor-oncology} \\
Source data & GSE176078 (HER2-enriched cells) \\
Key markers &
CDH3-high, AR-low, EZH2 $+118\%$ \\
Six Lock Type DOI &
\href{https://doi.org/10.5281/zenodo.18884158}
{10.5281/zenodo.18884158} \\
FOXA1/EZH2 Ratio DOI &
\href{https://doi.org/10.5281/zenodo.18883922}
{10.5281/zenodo.18883922} \\
AR Depth Axis DOI &
\href{https://doi.org/10.5281/zenodo.18891770}
{10.5281/zenodo.18891770} \\
Tazemetostat Fulvestrant DOI &
\href{https://doi.org/10.5281/zenodo.18884003}
{10.5281/zenodo.18884003} \\
License & CC BY 4.0 \\
\bottomrule
\end{tabular}

\end{document}
